% META: {"id":"TCOMPL-INEG-EX-004","chapitre":"TCOMPL-INEGALITES","type_objet":"exercice","capacites":["TCOMPL-INEG-C2"],"capacites_codes":["C2"],"methodes":[],"parcours":1,"competences":["raisonner"],"duree_min":10,"mode_creation":"ex_nihilo","sources_inspiration":[],"fichier_tex":"chapitres/TCOMPL-INEGALITES/exercices/TCOMPL-INEG-EX-004.tex","corrige_tex":"chapitres/TCOMPL-INEGALITES/corriges/TCOMPL-INEG-CO-004.tex","status":"approved","origin":"generated","status_history":["generated","audited","accepted","approved"],"release_acceptance":"RELEASE_OWNER_BATCH_ACCEPTANCE","acceptance_closure_digest":"sha256:9b3ccf9a81c5520fbb7e03b7b2d3e7bf2057a02834b6d908bf3be5f49f3b553f","acceptance_receipt_digest":"sha256:4fad86f0282e0fa4e0419cca1ad6379ae7af5a280c04a4a90880f90a1c044242"}
% BEGIN-VERIFY
% from sympy import Rational
% assert Rational(1, 2)**4 < Rational(1, 2)**2
% END-VERIFY
\begin{exercice}{TCOMPL-INEG-EX-004}{1}{10}
Deux pays sont modélisés par des courbes de Lorenz $L_A(x) = x^2$ et $L_B(x) = x^4$.
\begin{enumerate}
  \item Sans calculer d'intégrale, comparer $L_A(0{,}5)$ et $L_B(0{,}5)$ : que représentent ces deux valeurs concrètement ?
  \item En déduire, sans calcul, quel pays a la répartition la plus inégalitaire.
\end{enumerate}
\end{exercice}
